\documentclass[9pt,twocolumn]{extarticle}

\usepackage[letterpaper,margin=0.64in]{geometry}
\usepackage[T1]{fontenc}
\usepackage[utf8]{inputenc}
\usepackage{microtype}
\usepackage{lmodern}
\usepackage{amsmath,amssymb,amsthm,mathtools}
\usepackage{booktabs}
\usepackage{tabularx}
\usepackage{enumitem}
\usepackage{xcolor}
\usepackage[hidelinks]{hyperref}
\usepackage[capitalize,noabbrev]{cleveref}

\setlength{\columnsep}{0.22in}
\setlist{nosep,leftmargin=*}
\newtheorem{definition}{Definition}
\newtheorem{proposition}{Proposition}
\newtheorem{assumption}{Assumption}
\newcommand{\R}{\mathbb{R}}
\newcommand{\E}{\mathbb{E}}
\newcommand{\Pp}{\mathbb{P}}
\newcommand{\TV}{\operatorname{TV}}
\newcommand{\argmax}{\operatorname*{arg\,max}}
\newcommand{\ind}{\mathbf{1}}
\newcommand{\rtrue}{r^{\star}}
\newcommand{\rproxy}{\widetilde r}
\newcommand{\Dpre}{\mathcal{L}_{T}}
\newcommand{\PiFull}{\Pi_{\mathrm{full}}}
\newcommand{\PiHist}{\Pi_{\mathrm{hist}}}
\newcommand{\PiMem}{\Pi_{m}}

\title{Diagnosing Reward Misspecification and Observation Insufficiency\\
from Finite Prefixes of Reinforcement-Learning Training}
\author{%
Emilio Daza$^{1}$ \qquad Jacob Lee$^{2}$ \qquad Rohin Desai$^{3}$\\[0.45em]
\small $^{1}$\textit{Dartmouth College}\\
\small $^{2}$\textit{California Polytechnic State University}\\
\small $^{3}$\textit{UC Berkeley}}
\date{August 2026}

\begin{document}
\maketitle
\vspace{-1.2em}

\begin{abstract}
We ask when a finite prefix of ordinary reinforcement-learning (RL) training
identifies whether poor intended performance is caused by optimizing a
misspecified proxy reward or by lacking task-relevant observations. We define
the two failures as separate value gaps, distinguish intrinsic missing
information from limited agent memory, and formalize finite-prefix
identifiability as a property of distributions over passive logs. A two-action
construction proves that passive logs can be exactly observationally
equivalent under the two failure mechanisms. We then define privileged-state
and reward-audit probes and pose selection of a diagnostic intervention as a
minimum-cost statistical decision problem. The specification yields explicit
ground-truth labels and testable estimands for the accompanying experiments.
\end{abstract}

\section{Research question and scope}

\textbf{Research question.} When are reward misspecification and task-relevant
observation insufficiency statistically identifiable from a finite prefix of
RL training, and, when passive logs are insufficient, what minimum-cost
diagnostic intervention makes them identifiable?

Reward misspecification can allow greater proxy return while reducing intended
return \cite{pan2022effects}. Partial observability creates a different
difficulty: an agent must act without full knowledge of the controlled latent
state, and tractability requires structural conditions on how observations
reveal that state \cite{liu2022pomdp}. Early interactions are worth studying
because they can exert a lasting influence on deep-RL learning
\cite{nikishin2022primacy}. Our object is not merely prediction of eventual
failure, but identification of its causal class from early data.

\section{Formal setting}

Let a finite-horizon POMDP with two objectives be
\begin{equation}
\mathcal M=(\mathcal S,\mathcal A,\mathcal O,P,Z,\mu_0,H,
\gamma,\rproxy,\rtrue).
\end{equation}
Here $S_0\sim\mu_0$, $O_t\sim Z(\cdot\mid S_t)$,
$S_{t+1}\sim P(\cdot\mid S_t,A_t)$, $\rproxy:\mathcal
S\times\mathcal A\to[0,1]$ is the reward supplied during ordinary training,
and $\rtrue:\mathcal S\times\mathcal A\to[0,1]$ is the designer's intended
evaluation objective. The analysis also permits stochastic rewards by
interpreting these functions as conditional means. For objective $r$ and
policy $\pi$, define
\begin{equation}
V_r(\pi)=\E_{\mathcal M,\pi}\!\left[\sum_{t=0}^{H-1}
\gamma^t r(S_t,A_t)\right].
\end{equation}

Three policy classes isolate information from computation:
\begin{align*}
\PiFull &: \text{policies that observe the current latent state }S_t,\\
\PiHist &: \text{policies with perfect recall of the observable history},\\
\PiMem &: \text{policies restricted to a specified memory budget }m.
\end{align*}
The observable history before action $t$ is
$H_t^o=(O_0,A_0,\widetilde R_0,\ldots,O_t)$; thus $\PiHist$ may use all
information actually available to an ideal agent, whereas $\PiMem\subseteq
\PiHist$ represents the architecture being trained (for example, a reactive
policy or a fixed recurrent state).

For a learning algorithm $\mathsf A$, seed $U$, and $T$ training transitions,
the \emph{passive early-training prefix} is
\begin{equation}
\Dpre=(O_t,A_t,\widetilde R_t,p_t,V_t,\delta_t,d_t)_{t<T},U,
\end{equation}
where $p_t$ is the logged action distribution, $V_t$ a value estimate,
$\delta_t$ a TD error, and $d_t$ an episode boundary indicator. Latent states
$S_t$ and intended rewards $R_t^\star$ are deliberately excluded. Quantities
such as entropy or action diversity are measurable functions of $\Dpre$ and
therefore add no information in the statistical sense, although they may make
classification easier for a finite-capacity diagnostic model.

\section{Failure definitions}

For a policy class $\Pi$, let
$J_r(\Pi)=\sup_{\pi\in\Pi}V_r(\pi)$ and
$\Gamma_{\widetilde r}(\Pi)=\argmax_{\pi\in\Pi}V_{\widetilde r}(\pi)$.

\begin{definition}[Reward misspecification gap]
Holding the information interface fixed at $\PiHist$, define
\begin{equation}
\Delta_{\mathrm R}=
J_{\rtrue}(\PiHist)-
\sup_{\pi\in\Gamma_{\widetilde r}(\PiHist)}V_{\rtrue}(\pi).
\label{eq:reward-gap}
\end{equation}
Thus $\Delta_{\mathrm R}>0$ means that even the intended-best proxy-optimal
policy is suboptimal for the intended objective. For an actual learned policy
$\widehat\pi_T$, the empirical, algorithm-dependent regret is
$\widehat\Delta_{\mathrm R,T}=J_{\rtrue}(\PiHist)-
V_{\rtrue}(\widehat\pi_T)$ after controlling for the observation gap below.
\end{definition}

This definition is behavioral: $\rproxy\neq\rtrue$ alone is insufficient,
because two rewards may induce the same optimal policies. It also avoids
labeling an arbitrary unfavorable tie as structurally unavoidable.

\begin{definition}[Intrinsic observation gap]
\begin{equation}
\Delta_{\mathrm O}=J_{\rtrue}(\PiFull)-J_{\rtrue}(\PiHist).
\label{eq:observation-gap}
\end{equation}
If $\Delta_{\mathrm O}>0$, even an ideal policy with unlimited recall of the
available history cannot match a state-informed policy; task-relevant
information is intrinsically absent from the observation process.
\end{definition}

\begin{definition}[Memory/architecture gap]
\begin{equation}
\Delta_{\mathrm M}^{(m)}=J_{\rtrue}(\PiHist)-J_{\rtrue}(\PiMem).
\label{eq:memory-gap}
\end{equation}
A positive $\Delta_{\mathrm M}^{(m)}$ is not, by itself, observation
insufficiency: the required information exists in the observable history but
the chosen agent cannot retain or use it. This gap can be estimated by
comparing a history-aware oracle (or sufficiently expressive recurrent policy)
with the trained architecture.
\end{definition}

Given margins $\varepsilon_R,\varepsilon_O>0$, assign the population label
\begin{equation}
Y(\mathcal M)=
\big(\ind\{\Delta_{\mathrm R}>\varepsilon_R\},
     \ind\{\Delta_{\mathrm O}>\varepsilon_O\}\big)
\in\{0,1\}^2,
\end{equation}
corresponding to neither, reward-only, observation-only, and both. Margins
prevent negligible numerical differences from defining a failure.

\section{Finite-prefix identifiability}

Let $\Theta$ be a prespecified family of environment--objective pairs and let
$P_{\theta}^{T}$ denote the distribution of $\Dpre$ induced by
$\theta\in\Theta$, the fixed learner, and the seed distribution.

\begin{definition}[Passive identifiability at prefix length $T$]
The label is passively identifiable on $\Theta$ if
\begin{equation}
P_\theta^T=P_{\theta'}^T\quad\Longrightarrow\quad
Y(\theta)=Y(\theta')
\qquad\forall\theta,\theta'\in\Theta.
\label{eq:identifiability}
\end{equation}
Equivalently, every observational-equivalence class induced by the passive
prefix is label-homogeneous.
\end{definition}

For two simple hypotheses with equal prior probability, the optimal diagnostic
error after $n$ independent runs is
\begin{equation}
e_n^*=\frac{1-\TV\!\left((P_0^T)^{\otimes n},
(P_1^T)^{\otimes n}\right)}{2}.
\label{eq:lecam}
\end{equation}
Hence identical prefix distributions force chance-level binary diagnosis for
every sample size. Merely adding deterministic summaries of the same logs
cannot break this equivalence.

\subsection{An explicit passive impossibility result}

Fix $q\in(1/2,1)$ and consider independent one-step episodes with latent
$S\in\{0,1\}$, $\Pp(S=0)=q$, actions $A\in\{0,1\}$, and the constant
observation $O=o$. Define two systems:
\begin{align*}
\mathcal M_{\mathrm O}:\quad
&\widetilde R=R^\star=\ind\{A=S\};\\
\mathcal M_{\mathrm R}:\quad
&R^\star=\ind\{A=1\},\\[-0.25em]
&\widetilde R\mid A=a\sim
\begin{cases}
\mathrm{Bernoulli}(q),&a=0,\\
\mathrm{Bernoulli}(1-q),&a=1,
\end{cases}
\end{align*}
where $\widetilde R\perp S\mid A$ in $\mathcal M_{\mathrm R}$.

\begin{proposition}[Passive non-identifiability]
For any adaptive learner whose action at each episode is a measurable function
of its seed and previous passive logs,
\begin{equation}
P_{\mathcal M_{\mathrm O}}^T=P_{\mathcal M_{\mathrm R}}^T
\quad\text{for every finite }T.
\end{equation}
Nevertheless, $\mathcal M_{\mathrm O}$ is observation-limited but not
reward-misspecified, while $\mathcal M_{\mathrm R}$ is reward-misspecified but
not observation-limited.
\end{proposition}

\begin{proof}[Proof sketch]
In $\mathcal M_{\mathrm O}$, marginalizing the unobserved state gives
$\widetilde R\mid A=0\sim\mathrm{Bernoulli}(q)$ and
$\widetilde R\mid A=1\sim\mathrm{Bernoulli}(1-q)$, exactly as in
$\mathcal M_{\mathrm R}$. Because $O$ is constant, induction over episodes
shows equality of the full passive-log distributions under any adaptive action
rule. Logged policy probabilities, value estimates, and TD errors are produced
by the same learner from the same seed and history, so their distributions are
also equal.

In $\mathcal M_{\mathrm O}$, proxy and intended rewards coincide, so
$\Delta_{\mathrm R}=0$. A full-state policy chooses $A=S$ and obtains value
$1$, while the best history policy always chooses $A=0$ and obtains $q$;
therefore $\Delta_{\mathrm O}=1-q>0$. In $\mathcal M_{\mathrm R}$, the unique
proxy-optimal action is $0$, whereas the intended-optimal action is $1$, so
$\Delta_{\mathrm R}=1$. State access cannot improve the intended objective,
which is independent of $S$, and thus $\Delta_{\mathrm O}=0$.
\end{proof}

The construction also rules out explanations based only on reward scale: the
entire conditional distribution of the logged proxy reward is matched, not
just its mean.

\section{Diagnostic interventions}

An intervention $d$ modifies only the diagnostic data-collection protocol,
not the population label. We consider:
\begin{enumerate}[label=(\roman*)]
\item \textbf{Passive ($d_0$):} collect additional ordinary episodes.
\item \textbf{Privileged state ($d_S$):} reveal $S_t$ to a diagnostic probe for
selected decisions, while leaving the training reward unchanged.
\item \textbf{Reward audit ($d_R$):} query $R_t^\star$ on selected realized or
counterfactual state--action pairs, or run short probes with corrected reward.
\item \textbf{Combined ($d_{SR}$):} permit both forms of privileged access.
\end{enumerate}

The interventions target different causal contrasts. A state probe estimates
the gain $J_{\rtrue}(\PiFull)-J_{\rtrue}(\PiHist)$; a reward audit estimates how
the intended ranking of actions differs from the proxy ranking. In the toy
pair, privileged state predicts proxy reward in $\mathcal M_{\mathrm O}$ but is
irrelevant to proxy reward in $\mathcal M_{\mathrm R}$, whereas an audit finds
$R^\star=\widetilde R$ in $\mathcal M_{\mathrm O}$ but not in
$\mathcal M_{\mathrm R}$. Either targeted probe therefore breaks the passive
equivalence statistically; using both is needed for unrestricted four-class
diagnosis when the two failures may coexist.

Let $N_d$ be probe episodes, $Q_S(d)$ privileged-state queries, and $Q_R(d)$
intended-reward queries. For application-dependent nonnegative prices, define
\begin{equation}
C(d)=c_E N_d+c_S Q_S(d)+c_R Q_R(d).
\label{eq:cost}
\end{equation}
Writing $P_\theta^{T,d}$ for all passive and interventional data, the
minimum-cost $\delta$-reliable intervention is
\begin{equation}
C_\delta^*=\inf_{d,g}\left\{C(d):
\sup_{\theta\in\Theta}
\Pp_\theta\!\left[g(\Dpre,D_d)\neq Y(\theta)\right]\leq\delta
\right\}.
\label{eq:min-cost}
\end{equation}
Here $g$ ranges over diagnostic rules and $D_d$ is the probe data. A necessary
pairwise condition is
\begin{equation}
\inf_{Y(\theta)\neq Y(\theta')}
\TV(P_\theta^{T,d},P_{\theta'}^{T,d})\geq 1-2\delta.
\end{equation}
Thus the ``minimum'' intervention is not universally one query; it depends on
the candidate family, noise level, confidence target, and relative cost of
state and reward access.

\section{Assumptions for positive identification}

The impossibility result shows that identification requires restrictions or
interventions. The initial experiments should state and test the following:
\begin{enumerate}[label=\textbf{A\arabic*.}]
\item \textbf{Known protocol.} The learner, logging map, prefix length, and seed
distribution are fixed or observed; train and diagnostic trajectories do not
silently change across conditions.
\item \textbf{Exploration/overlap.} Every action needed to distinguish the
hypotheses is selected with probability at least $\eta>0$ in relevant
histories, or is deliberately selected by a probe.
\item \textbf{Margin and stationarity.} Failure gaps exceed their declared
thresholds, and transition, observation, and reward mechanisms remain stable
during a run.
\item \textbf{Faithful probes.} Privileged state reveals task-relevant state
without changing dynamics; the reward audit is an unbiased measurement of the
declared intended objective.
\item \textbf{Structural separation.} Within the chosen environment family,
every cross-label pair is separated by the selected passive or interventional
data distribution. Without this assumption, no classifier can guarantee
causal identification.
\end{enumerate}

Under A1--A5 and independent runs, standard consistent hypothesis tests recover
the label as the number of runs grows. Finite-sample claims should be reported
with uncertainty rather than inferred from a small set of point estimates;
this is especially important in few-run RL evaluation
\cite{agarwal2021precipice}.

\section{Experimental estimands and handoff}

The environment code should store $S_t$, $R_t^\star$, and $\widetilde R_t$
separately, exposing only the latter in passive logs. The following oracle
comparisons validate the ground-truth condition before training a diagnostic:
\begin{center}
\scriptsize
\begin{tabularx}{\columnwidth}{@{}lXX@{}}
\toprule
Quantity & Population contrast & Interpretation\\
\midrule
$\Delta_{\mathrm R}$ & intended-best vs. proxy-optimal & incentive failure\\
$\Delta_{\mathrm O}$ & full-state vs. full-history & missing information\\
$\Delta_{\mathrm M}^{(m)}$ & full-history vs. memory-$m$ & agent limitation\\
\bottomrule
\end{tabularx}
\end{center}

Use common random numbers or paired seeds across the four conditions where
possible. Passive classifiers receive only $\Dpre$ at the first $5\%$, $10\%$,
and $20\%$ of training; privileged labels never enter passive features. Report
balanced accuracy, macro-F1, per-class recall, confusion matrices, and
uncertainty intervals across seeds. The central empirical curve is diagnostic
accuracy versus both prefix length $T$ and intervention cost $C(d)$. Five paired
seeds are suitable for pipeline validation, but not by themselves for precise
claims; final run counts should be chosen from pilot variance and accompanied
by interval estimates.

\section{Limitations and falsifiable claims}

The labels depend on a declared intended reward; no statistical procedure can
recover an unspecified human objective. Oracle policies may be approximated in
larger environments, confounding optimization error with the value gaps.
Moreover, diagnosis is relative to $\Theta$: unrestricted POMDP families admit
observationally equivalent explanations. The project should therefore test
three falsifiable claims: (i) passive accuracy remains near chance on matched
equivalence pairs; (ii) the appropriate single probe separates each pure
failure class; and (iii) combined probes achieve reliable four-class diagnosis
at lower cost than indiscriminate privileged data collection.

\section{Conclusion}

Reward misspecification and observation insufficiency are distinct population
properties, but they need not be distinguishable from ordinary early-training
logs. The correct theoretical baseline is therefore conditional: identify the
label only under explicit structural assumptions, and otherwise choose the
least costly intervention that separates cross-label data distributions.

\begin{thebibliography}{9}
\small

\bibitem{pan2022effects}
A. Pan, K. Bhatia, and J. Steinhardt.
\newblock The effects of reward misspecification: Mapping and mitigating
misaligned models.
\newblock In \emph{International Conference on Learning Representations},
2022.

\bibitem{liu2022pomdp}
Q. Liu, A. Chung, C. Szepesv\'ari, and C. Jin.
\newblock When is partially observable reinforcement learning not scary?
\newblock In \emph{Proceedings of the 35th Conference on Learning Theory},
PMLR 178:5175--5220, 2022.

\bibitem{nikishin2022primacy}
E. Nikishin, M. Schwarzer, P. D'Oro, P.-L. Bacon, and A. Courville.
\newblock The primacy bias in deep reinforcement learning.
\newblock In \emph{Proceedings of the 39th International Conference on Machine
Learning}, PMLR 162:16828--16847, 2022.

\bibitem{agarwal2021precipice}
R. Agarwal, M. Schwarzer, P. S. Castro, A. C. Courville, and M. G. Bellemare.
\newblock Deep reinforcement learning at the edge of the statistical
precipice.
\newblock In \emph{Advances in Neural Information Processing Systems},
volume 34, 2021.

\end{thebibliography}

\end{document}
